\documentclass[11pt,a4paper]{article}

\usepackage[T1]{fontenc}
\usepackage[utf8]{inputenc}
\usepackage{lmodern}
\usepackage[margin=2.2cm]{geometry}
\usepackage{amsmath,amssymb}
\usepackage{booktabs}
\usepackage{array}
\usepackage{xcolor}
\usepackage{enumitem}
\usepackage{fancyhdr}
\usepackage[hidelinks]{hyperref}

\definecolor{ink}{HTML}{222222}
\definecolor{muted}{HTML}{5B5B5B}
\definecolor{rule}{HTML}{D8D8D8}
\definecolor{boxbg}{HTML}{F7F7F5}
\definecolor{seismic}{HTML}{3573B9}
\definecolor{electric}{HTML}{C98A00}
\definecolor{magnet}{HTML}{B03A5B}

\pagestyle{fancy}
\fancyhf{}
\lhead{\small\color{muted}SciML Project --- the multi-physics idea, explained simply}
\rhead{\small\color{muted}\thepage}
\renewcommand{\headrulewidth}{0.4pt}

\setlength{\parskip}{4pt}
\setlength{\parindent}{0pt}
\emergencystretch=3em

% A soft grey box for little stories and for the "grown-up corner".
\newsavebox{\storyboxsave}
\newenvironment{storybox}[1]
  {\par\medskip\noindent
   \begin{lrbox}{\storyboxsave}
   \begin{minipage}{\dimexpr\linewidth-2\fboxsep-2\fboxrule\relax}
   \textbf{#1}\par\smallskip}
  {\end{minipage}\end{lrbox}%
   \fcolorbox{rule}{boxbg}{\usebox{\storyboxsave}}\par\medskip}

\title{\textbf{Looking Inside the Earth Without Digging}\\[4pt]
\large The ``multi-physics dataset'' idea, explained like you are five}
\author{SciML Project}
\date{September 2026}

\begin{document}
\maketitle

\section*{The problem: a cake you are not allowed to cut}

Imagine the ground under your feet is a giant layered cake.
There is a layer of mud, then a layer of sand, then a layer of hard rock,
then maybe a layer of wet, salty sponge, and so on, all the way down.

We really want to know what the layers are and where they are.
People who look for water, oil, metals, or safe places to build all need this.
But the cake is \emph{enormous} and we are not allowed to cut it.
Digging is slow, expensive, and only tells you about one tiny spot.

So instead of cutting the cake, we \textbf{poke it and listen}.
Every way of poking the cake tells us something different.
This whole project is about three ways of poking, and about teaching a computer
to turn the ``listening'' back into a picture of the cake.

\section*{Three ways to poke the ground}

\subsection*{1. \textcolor{seismic}{Thump it} (sound)}

Stomp on the ground. A wobble travels down into the cake.
When the wobble hits the edge between two layers (say, soft sand on top of hard rock),
part of it bounces back up. We put little ``ears'' (called geophones) on the surface
and record the bounces coming back.

The wobble travels \emph{fast} through hard rock and \emph{slow} through soft mud.
So if we know how fast the wobble goes everywhere, we know what kind of stuff is where.

\begin{itemize}[leftmargin=2em]
  \item \textbf{What the ground \emph{is}:} a map of \textbf{how fast sound goes} in every spot.
        Grown-ups call it \textbf{velocity}, letter \textbf{v}.
  \item \textbf{What we \emph{hear}:} the wiggles recorded by the ears over time.
        Grown-ups call it the \textbf{pressure} record or seismogram, letter \textbf{p}.
\end{itemize}

\begin{storybox}{Picture it as: a bouncing ball}
You throw a ball down a staircase. It bounces off every step. From the timing of the
bounces you can work out where the steps are. The ball keeps its ``shape'' as it
bounces --- this is a \emph{wave}.
\end{storybox}

\subsection*{2. \textcolor{electric}{Push electricity into it} (current)}

Stick two metal rods into the ground and push electricity in through one rod and
out through the other. Some parts of the cake let electricity flow easily
(wet clay, salty water). Some parts fight it (dry rock, sand, ice).
We stick more rods around and measure how much electricity arrives where.

\begin{itemize}[leftmargin=2em]
  \item \textbf{What the ground \emph{is}:} a map of \textbf{how hard it is to push electricity through} each spot.
        Grown-ups call it \textbf{resistivity}, letter \textbf{r}.
  \item \textbf{What we \emph{measure}:} how much \textbf{current} flows and what voltage shows up at the rods,
        letter \textbf{i}.
\end{itemize}

\begin{storybox}{Picture it as: water finding its level in a pond}
Pour water into one corner of a pond full of rocks and sponges.
It spreads out and settles. Nothing bounces, nothing wobbles; it just finds a
resting shape. Where the water sits tells you where the sponges are.
This is a \emph{steady} problem: time does not matter, only the final shape.
\end{storybox}

\subsection*{3. \textcolor{magnet}{Wiggle a magnet over it} (magnetism)}

Hold a big magnet (or a big coil of wire) above the ground and wiggle it back and forth.
A wiggling magnet makes electricity swirl around inside anything that can conduct ---
like little whirlpools of electricity in the wet, salty layers.
Those whirlpools make \emph{their own} tiny magnet field, which comes back up
and we can feel it at the surface with a sensor.

Layers that conduct electricity well make big whirlpools and a strong answer.
Layers that do not conduct stay quiet.

\begin{itemize}[leftmargin=2em]
  \item \textbf{What the ground \emph{is}:} a map of \textbf{how easily electricity flows} in each spot.
        Grown-ups call it \textbf{conductivity}, letter \textbf{c}.
  \item \textbf{What we \emph{feel}:} the \textbf{magnetic} field that comes back, letter \textbf{m}.
\end{itemize}

\begin{storybox}{Picture it as: a drop of ink in water}
Drop ink into a glass of water. It does not bounce; it slowly smears outward
and fades. The wiggling-magnet signal behaves like that inside the ground:
it \emph{spreads and fades} instead of bouncing. This is a \emph{diffusion} problem.
The rules for it are called \textbf{Maxwell's equations}, and ``in the geophysical
perspective'' just means ``the slow, smeary version that works inside the earth''.
\end{storybox}

\begin{storybox}{Wait --- are r and c the same thing?}
Yes! ``Hard to push electricity through'' and ``easy to push electricity through''
are the same fact said two ways, like ``heavy'' and ``light''.
Grown-ups write $r = 1/c$.
So why list both? Because they belong to two \emph{different games}:
the push-electricity game (steady, like the pond) and the wiggle-magnet game
(smeary, like the ink). The map is the same; the poking is different,
so what we measure is different.
\end{storybox}

\section*{Pairs: ``what the cake is'' and ``what we heard''}

Every game gives us a \textbf{pair}: a picture of the cake, and the recording we got by poking it.

\begin{center}
\begin{tabular}{@{}lllll@{}}
\toprule
Game & Poke & What the cake is & What we record & The pair \\
\midrule
\textcolor{seismic}{Thump}        & stomp        & velocity (v)     & wiggles / pressure (p) & (v, p) \\
\textcolor{magnet}{Wiggle magnet} & shake a coil & conductivity (c) & magnetic field (m)     & (c, m) \\
\textcolor{electric}{Push current}& metal rods   & resistivity (r)  & current / voltage (i)  & (r, i) \\
\bottomrule
\end{tabular}
\end{center}

When the message said ``\emph{v,p pair --- I want c,m and r,i}'' it meant exactly this:
\emph{we already have the thump pair; now make the magnet pair and the electricity pair too.}

\subsection*{The game the computer plays}

Going \emph{forward} is the easy direction: if you know the cake, a physics program can
work out what the ears would hear. That is the ``poke and listen'' simulation.

The direction we actually care about is \emph{backwards}: we only have the recording,
and we want to draw the cake. That is hard --- like guessing what is inside a wrapped
present by shaking it.

So we play a practice game with the computer:
\begin{enumerate}[leftmargin=2em]
  \item We invent thousands of pretend cakes (we know exactly what is inside each one).
  \item We run the physics program to get the recording for each pretend cake.
  \item We show the computer the recording and ask it to draw the cake.
  \item We show it the real answer, it fixes its mistakes, and we repeat --- thousands of times.
\end{enumerate}
After enough practice, we hand it a \emph{real} recording and it draws a cake we have never seen.
This is what all the neural networks in the project (FNO, DeepONet, PINNs, KANs) are doing.

\section*{Why use all three games instead of one?}

Each way of poking is good at different things:

\begin{itemize}[leftmargin=2em]
  \item \textcolor{seismic}{\textbf{Thumping}} is great at finding \emph{edges} between hard and soft layers.
        But it cannot tell wet sand from dry sand very well.
  \item \textcolor{electric}{\textbf{Pushing electricity}} is great at finding \emph{wet and salty} stuff.
        But it gets blurry very quickly the deeper you go.
  \item \textcolor{magnet}{\textbf{Wiggling a magnet}} can reach \emph{deep} conductors that the rods cannot.
        But it is smeary and cannot see sharp edges.
\end{itemize}

\begin{storybox}{Picture it as: guessing a fruit with your eyes closed}
Touch alone might say ``round''. Smell alone might say ``sweet''. Taste alone might say ``sour''.
Put them together and you know it is an orange. Using all the clues together has a
fancy name: \textbf{joint inversion}.
\end{storybox}

There is one more trick. All three games are poking the \emph{same cake}.
So wherever the sound map has a layer edge, the electricity map should have an edge
in the same place too. The computer is allowed to use that rule (``the pictures must agree
on where the edges are'') to fix each picture using the others.

\section*{The dataset: a big box of pretend cakes}

This is the ``build a dataset with many physical parameters'' part of the message.
Today we have a box of pretend cakes that only know their \emph{sound speed}.
We want a box where every cake knows \textbf{all three} things at once:

\begin{center}
\begin{tabular}{@{}ll@{}}
\toprule
For every pretend cake we store\ldots & \\
\midrule
the layers themselves (the ``recipe'') & mud here, sand there, rock below \\
velocity map (v) $\to$ thump recording (p) & the old game \\
conductivity map (c) $\to$ magnet recording (m) & new \\
resistivity map (r) $\to$ current recording (i) & new \\
\bottomrule
\end{tabular}
\end{center}

The important rule is that the three maps must come from \emph{one} recipe.
Real rocks link these things: a wet, spongy layer is \emph{slow} for sound \emph{and}
\emph{easy} for electricity; hard dry rock is \emph{fast} for sound \emph{and} \emph{hard}
for electricity. Grown-ups have little formulas for these links (with names like Archie,
Gardner, Wyllie, Faust). We draw the recipe once, then use the formulas to paint
v, c and r from it. That way the three pictures are honest about being the same cake.

\section*{What we already have, and what is missing}

\textbf{Already built:}
\begin{itemize}[leftmargin=2em]
  \item The \textcolor{seismic}{thump game}, in two flavours: our own 1D cake (the PINN/KAN experiments),
        and the big public 2D cake collection called OpenFWI, plus the neural networks that learned it.
  \item The physics programs that simulate thumping (finite differences, and friends).
\end{itemize}

\textbf{Missing:}
\begin{itemize}[leftmargin=2em]
  \item The recipe tool: ``draw one cake, paint v, c and r from it''.
  \item The \textcolor{electric}{push-electricity} simulator (the pond).
  \item The \textcolor{magnet}{wiggle-magnet} simulator (the ink).
  \item Neural networks that play the new games, and one that plays all three together.
\end{itemize}

\textbf{What to do first:} the \textcolor{electric}{push-electricity game}.
Its physics is the simplest of the three (no time, nothing moves --- just the pond finding its level),
it gives us a second pair on the very same cakes right away, and it tests the whole
pipeline --- recipe $\to$ simulator $\to$ dataset $\to$ network --- before we tackle the smeary magnet game.

\section*{Grown-up corner (you can skip this)}

The three games are three different \emph{kinds} of maths problem, which is a nice reason
to have them all in one benchmark:

\begin{center}
\begin{tabular}{@{}llll@{}}
\toprule
Game & Kind of equation & Picture & Equation (2D) \\
\midrule
\textcolor{seismic}{Thump} & hyperbolic (wave) & bouncing ball &
  $\rho\, u_{tt} = \nabla\!\cdot\!\big(E\,\nabla u\big)$ \\
\textcolor{magnet}{Wiggle magnet} & parabolic (diffusion) & ink in water &
  $\nabla^2 \mathbf{E} = i\omega\mu\sigma\,\mathbf{E}$ \quad (quasi-static Maxwell) \\
\textcolor{electric}{Push current} & elliptic (steady) & pond level &
  $\nabla\!\cdot\!\big(\sigma\,\nabla\phi\big) = -I\,\delta(\mathbf{x}-\mathbf{x}_s)$ \\
\bottomrule
\end{tabular}
\end{center}

\begin{itemize}[leftmargin=2em]
  \item ``Geophysical perspective'' on Maxwell means dropping the displacement current
        (the earth is a conductor, frequencies are low), which turns Maxwell into a
        diffusion equation. How deep the magnet game can see is set by the
        \emph{skin depth}, roughly $\delta \approx 503\sqrt{\rho/f}$ metres.
  \item The push-current game is the same equation with $\omega \to 0$: pure Poisson.
  \item The ``pictures must agree on edges'' rule is the \emph{cross-gradient} constraint
        $\nabla v \times \nabla \sigma = 0$ used in joint inversion.
  \item Same geology, three PDE classes, one operator-learning benchmark --- that is the
        one-sentence pitch.
\end{itemize}

\section*{The whole thing in one breath}

The ground is a cake we cannot cut. We can thump it, push electricity into it, or
wiggle a magnet over it, and each gives us a different recording. We already taught a
computer the thump game. Now we want to build a box of pretend cakes that know all
three properties at once, simulate all three recordings for each, and teach the
computer to draw the cake from any of them --- or, best of all, from all three together.

\end{document}
